% =========================================================
% Reference Page: Axiom Systems for N and Z
% File: volume-ii/integers/notes/notes-nz-axioms-reference.tex
% =========================================================

\section{Reference: Axiom Systems for $\mathbb{N}$ and $\mathbb{Z}$}
\label{sec:nz-axioms-reference}

% ---------------------------------------------------------
% PEANO AXIOMS FOR N
% ---------------------------------------------------------

\begin{tcolorbox}[colback=axiombox, colframe=axiomborder, arc=2pt,
  left=6pt, right=6pt, top=4pt, bottom=4pt,
  title={\small\textbf{Peano Axioms for $\mathbb{N}$ (Tao \S2.1)}},
  fonttitle=\small\bfseries]

Let $\mathbb{N}$ be a set containing a distinguished element $0$
and equipped with a \emph{successor} (increment) operation
$n \mapsto n{+\!\!+}$.

\medskip
\begin{tabular}{@{} l p{0.80\linewidth} @{}}
\textbf{P1} & $0$ is a natural number. \\[4pt]
\textbf{P2} & If $n \in \mathbb{N}$, then $n{+\!\!+} \in \mathbb{N}$.
             \hfill\textit{(closure under successor)} \\[4pt]
\textbf{P3} & $n{+\!\!+} \neq 0$ for all $n \in \mathbb{N}$.
             \hfill\textit{(zero is not a successor)} \\[4pt]
\textbf{P4} & If $n{+\!\!+} = m{+\!\!+}$, then $n = m$.
             \hfill\textit{(successor is injective)} \\[4pt]
\textbf{P5} & Let $P$ be any property of natural numbers.
             If $P(0)$ holds, and whenever $P(n)$ holds we also have
             $P(n{+\!\!+})$, then $P(n)$ holds for all $n \in \mathbb{N}$.
             \hfill\textit{(principle of mathematical induction)} \\
\end{tabular}
\end{tcolorbox}

\begin{remark}[What the Peano axioms do and do not give]
P1--P4 rule out three pathologies: $\mathbb{N}$ wrapping around
(e.g.\ $3{+\!\!+} = 0$), the successor hitting a ceiling
(e.g.\ $4{+\!\!+} = 4$), and rogue elements unreachable from $0$.
P5 eliminates the last: the only elements of $\mathbb{N}$ are those
reachable by starting at $0$ and applying the successor operation
finitely many times.

The Peano axioms say nothing about addition, multiplication, or order.
These are \emph{defined} recursively from the successor and then
shown to satisfy familiar algebraic laws by induction.
\end{remark}

\begin{remark}[Equivalence with well-ordering]
The following are logically equivalent over $\mathbb{N}$:
\begin{itemize}
  \item The Principle of Mathematical Induction (P5).
  \item The Well-Ordering Principle: every nonempty subset of
        $\mathbb{N}$ has a least element.
  \item Strong induction: if $P(k)$ holds for all $k < n$ implies
        $P(n)$, then $P$ holds everywhere.
\end{itemize}
One must assume at least one of these; the others then follow.
\end{remark}

\medskip

% ---------------------------------------------------------
% RING AXIOMS FOR Z
% ---------------------------------------------------------

\begin{tcolorbox}[colback=axiombox, colframe=axiomborder, arc=2pt,
  left=6pt, right=6pt, top=4pt, bottom=4pt,
  title={\small\textbf{Axioms for $\mathbb{Z}$ as an Ordered Commutative Ring}},
  fonttitle=\small\bfseries]

The integers $\mathbb{Z}$ are equipped with addition $+$,
multiplication $\cdot$, and order $\leq$.

\medskip
\noindent\textbf{Addition axioms}

\medskip
\begin{tabular}{@{} l p{0.78\linewidth} @{}}
\textbf{A1} & $a + b = b + a$
             \hfill\textit{(commutativity of addition)} \\[4pt]
\textbf{A2} & $(a + b) + c = a + (b + c)$
             \hfill\textit{(associativity of addition)} \\[4pt]
\textbf{A3} & There exists $0 \in \mathbb{Z}$ such that $a + 0 = a$.
             \hfill\textit{(additive identity)} \\[4pt]
\textbf{A4} & For every $a \in \mathbb{Z}$, there exists $-a \in \mathbb{Z}$
             such that $a + (-a) = 0$.
             \hfill\textit{(additive inverses)} \\[4pt]
\end{tabular}

\medskip
\noindent\textbf{Multiplication axioms}

\medskip
\begin{tabular}{@{} l p{0.78\linewidth} @{}}
\textbf{M1} & $a \cdot b = b \cdot a$
             \hfill\textit{(commutativity of multiplication)} \\[4pt]
\textbf{M2} & $(a \cdot b) \cdot c = a \cdot (b \cdot c)$
             \hfill\textit{(associativity of multiplication)} \\[4pt]
\textbf{M3} & There exists $1 \in \mathbb{Z}$, $1 \neq 0$, such that
             $a \cdot 1 = a$.
             \hfill\textit{(multiplicative identity)} \\[4pt]
\end{tabular}

\medskip
\noindent\textbf{Distributive law}

\medskip
\begin{tabular}{@{} l p{0.78\linewidth} @{}}
\textbf{D} & $a \cdot (b + c) = a \cdot b + a \cdot c$
            \hfill\textit{(distributivity)} \\[4pt]
\end{tabular}

\medskip
\noindent\textbf{Order axioms}

\medskip
\begin{tabular}{@{} l p{0.78\linewidth} @{}}
\textbf{O1} & For all $a, b$: $a \leq b$ or $b \leq a$.
             \hfill\textit{(totality)} \\[4pt]
\textbf{O2} & If $a \leq b$ and $b \leq a$, then $a = b$.
             \hfill\textit{(antisymmetry)} \\[4pt]
\textbf{O3} & If $a \leq b$ and $b \leq c$, then $a \leq c$.
             \hfill\textit{(transitivity)} \\[4pt]
\textbf{O4} & If $b \leq c$, then $a + b \leq a + c$.
             \hfill\textit{(order compatible with addition)} \\[4pt]
\textbf{O5} & If $a \geq 0$ and $b \geq 0$, then $ab \geq 0$.
             \hfill\textit{(order compatible with multiplication)} \\[4pt]
\end{tabular}

\end{tcolorbox}

\begin{remark}[What $\mathbb{Z}$ gains over $\mathbb{N}$]
The decisive addition is \textbf{A4}: additive inverses.
$\mathbb{N}$ satisfies A1--A3, M1--M3, D, and O1--O5,
but it has no additive inverses ($3 - 5 \notin \mathbb{N}$).
$\mathbb{Z}$ is obtained from $\mathbb{N}$ by formally adjoining
a negative $-n$ for each $n$, making subtraction always defined.
\end{remark}

\begin{remark}[What $\mathbb{Z}$ lacks: no multiplicative inverses]
$\mathbb{Z}$ satisfies all the field axioms \emph{except} for
the multiplicative inverse axiom.
There is no $x \in \mathbb{Z}$ with $2x = 1$, so M4
(``every nonzero element has a multiplicative inverse'') fails.
In the language of abstract algebra, $\mathbb{Z}$ is a
\emph{commutative ordered ring} but not a field.
The passage to $\mathbb{Q}$ restores M4 by formally adjoining
fractions $p/q$.
\end{remark}

\medskip

% ---------------------------------------------------------
% COMPARISON TABLE
% ---------------------------------------------------------

\begin{tcolorbox}[colback=gray!6, colframe=gray!40, arc=2pt,
  left=6pt, right=6pt, top=4pt, bottom=4pt,
  title={\small\textbf{Axiom Comparison: $\mathbb{N}$, $\mathbb{Z}$,
         $\mathbb{Q}$, $\mathbb{R}$}},
  fonttitle=\small\bfseries]

\small
\renewcommand{\arraystretch}{1.3}
\begin{tabular}{@{} l l c c c c @{}}
\toprule
\textbf{Group} & \textbf{Axiom} &
$\mathbb{N}$ & $\mathbb{Z}$ & $\mathbb{Q}$ & $\mathbb{R}$ \\
\midrule
Addition    & A1 commutativity          & \checkmark & \checkmark & \checkmark & \checkmark \\
            & A2 associativity          & \checkmark & \checkmark & \checkmark & \checkmark \\
            & A3 identity ($0$)         & \checkmark & \checkmark & \checkmark & \checkmark \\
            & A4 inverses ($-a$)        & $\times$   & \checkmark & \checkmark & \checkmark \\
\midrule
Multiplication & M1 commutativity       & \checkmark & \checkmark & \checkmark & \checkmark \\
            & M2 associativity          & \checkmark & \checkmark & \checkmark & \checkmark \\
            & M3 identity ($1$)         & \checkmark & \checkmark & \checkmark & \checkmark \\
            & M4 inverses ($1/a$)       & $\times$   & $\times$   & \checkmark & \checkmark \\
\midrule
            & D distributivity          & \checkmark & \checkmark & \checkmark & \checkmark \\
\midrule
Order       & O1--O3 total order        & \checkmark & \checkmark & \checkmark & \checkmark \\
            & O4--O5 order + arithmetic & \checkmark & \checkmark & \checkmark & \checkmark \\
\midrule
Completeness & Least upper bound        & $\times$   & $\times$   & $\times$   & \checkmark \\
\midrule
\multicolumn{2}{@{}l}{\textbf{Structure}} &
\textit{semiring} & \textit{ring} & \textit{field} &
\textit{complete} \\
\multicolumn{2}{@{}l}{} & & & & \textit{ordered field} \\
\bottomrule
\end{tabular}

\medskip
\noindent\small
$\checkmark$ = satisfied \quad $\times$ = fails
\end{tcolorbox}

\begin{remark}[The hierarchy of number systems]
Each extension repairs exactly one deficiency:
\[
\mathbb{N}
\;\xrightarrow{+\;\text{A4}}\;
\mathbb{Z}
\;\xrightarrow{+\;\text{M4}}\;
\mathbb{Q}
\;\xrightarrow{+\;\text{LUB}}\;
\mathbb{R}.
\]
$\mathbb{N} \to \mathbb{Z}$: adjoin additive inverses so subtraction is
always defined.
$\mathbb{Z} \to \mathbb{Q}$: adjoin multiplicative inverses so division
by nonzero elements is always defined.
$\mathbb{Q} \to \mathbb{R}$: adjoin least upper bounds so that limits of
Cauchy sequences always converge.
Each step is forced by a specific incompleteness of the previous system.
\end{remark}

\begin{remark}[Well-ordering as the characteristic property of $\mathbb{N}$]
No row in the table captures the property that most distinguishes
$\mathbb{N}$ from $\mathbb{Z}$: the Well-Ordering Principle.
$\mathbb{N}$ is well-ordered (every nonempty subset has a least element);
$\mathbb{Z}$ is not (the set of all negative integers has no least element).
Well-ordering is the engine behind the Division Algorithm, Bézout's
Identity, and ultimately unique prime factorization.
\end{remark}